\begin{summarybox}
	\textbf{\textcolor{CSummary}{一句话核心}}

	\textbf{若 $\xi\sim B(n,p)$，则二项分布的数学期望为 $E\xi=np$。其中 $n$ 是试验次数，$p$ 是每次试验的成功概率。}

	\medskip
	\textbf{\textcolor{CSummary}{使用条件}}
	\begin{itemize}[leftmargin=*, itemsep=0.5em, topsep=0.4em]
		\item \textbf{直接套用二项分布公式：} 需要先确认 $\xi\sim B(n,p)$，即 $n$ 次独立重复试验，每次成功概率均为 $p$。
		\item \textbf{更一般的求期望方法：} 若 $\xi$ 表示成功次数，可写成指示变量之和
		      \[
			      \xi=X_1+X_2+\cdots+X_n,
		      \]
		      则
		      \[
			      E\xi=\sum_{i=1}^n EX_i.
		      \]
		      若各次成功概率均为 $p$，则仍有 $E\xi=np$。
	\end{itemize}

	\medskip
	\textbf{\textcolor{CSummary}{关键提醒}}
	\begin{itemize}[leftmargin=*, itemsep=0.5em, topsep=0.4em]
		\item \textbf{公式区分：} 二项分布中 $E\xi=np$，方差为 $D\xi=np(1-p)$。
		\item \textbf{模型验证：} 判断分布时必须检查独立重复；但单独求期望时，可以优先考虑指示变量法。
	\end{itemize}
\end{summarybox}